% !TeX root = ../papers/nc-p3.tex
\section{両側が楕円多重点となる被約極限}
\label{sec:reduced-additive-limits}
\begin{theorem}[被約二側の特殊化]\label{thm:all-reduced-apolar-specializations}
滑らかな楕円鉛筆の特殊化で、両方の極限二次式の鉛筆が被約な四次完全交叉を切るとする。
特殊反転窓が標準ASなら、両曲線は同じ被約曲線として同定でき、二つの埋込み束は等しい
正の多重次数を持ち、六関係はその上の基点自由鉛筆から得られる。従って反対階数は6または8である。
\end{theorem}
\begin{proof}
整な側を持つ場合と節点多角形の場合は前二節による。残る両曲線$X,Y$は唯一の楕円多重点
$x_*,y_*$を持ち、いずれも二本以上の滑らかな有理枝からなる。各枝の次数を$d_i,e_j$、
枝数を$p,q$とする。総次数は各々4で、全て正である。埋込み束の切断空間を$A,B$とする。
各枝への制限は完全線型系である。他枝で一次係数の和を調整できるためである。
特異点で消え、指定一枝だけに支えを持つ切断はその原点で二重に消えるもの全体である。

一般点の同定のグラフを二曲線族の積内で閉包し、特殊ファイバーの被約支持を$G$とする。
両射影は正規な全空間への固有双有理射なのでファイバーは連結である。
被約な特殊曲線をサイクルとして押し出す議論は前節と同じで、各枝を支配する成分はただ一つで
その次数は1である（押し出すのはグラフの特殊ファイバーの重複度付きサイクルである）。
従って$G$の成分は、二つの枝の同型のグラフ、水平枝$X_i\times\{y_i\}$、
垂直枝$\{x_j\}\times Y_j$である。前者を対応枝と呼ぶ。
Segre埋込みで$G$の点評価が張るベクトル空間の次元$r$は積$AB$の被約成分への
制限の階数であり、補題\ref{lem:mixed-rank-loss}により$r\le8$である。

以下で連結な成分の合併のSegre線型包を使う。その両射影のベクトル空間次元が$a,b$なら、
この包のベクトル空間次元は$a+b-1$以上である。実際、一成分の二射影の切断は
射影直線の多項式で、その積の空間は次元の和から1を引いた値以上である
（互いに異なる先頭次数の基底をとり、先頭次数の和集合を見る）。
既存の連結な部分にそれと交わる一成分を加え、両側の増分を$c,d$とする。
$c>0$なら既存部分で消える左切断の新成分への像は$c$次元で、
右切断の制限は$d$に交点での非零評価を足した$d+1$次元以上である。
積は既存部分で消え、増分は$c+d$以上となる。$c=0,d>0$では左右を交換する。
成分を順に加えれば主張を得る。特に$7\le r\le8$である。
$\Phi$を$G$上の積とすれば、$\Phi_x$の階数が1以下の非零文字$x$は許されない。
実際の評価は$r=7$でも階数を高々1しか増やせず、二次元の評価核が指数成長を強制する。
$r=8$なら$\Phi_x$の階数は各非零文字で3以上必要である。

\paragraph{二つの特異点の対がグラフの外にある場合。}
$(x_*,y_*)\notin G$とする。$x_*$のファイバーは正次数の垂直枝を含めず連結なので一点$y_0\ne y_*$、
同様に$y_*$上は一点$x_0\ne x_*$である。全水平枝は$y_0$に、全垂直枝は$x_0$に載る。
対応枝は$x_0$の枝と$y_0$の枝の間の高々一つである。
対応枝がなければ$x_0,y_0$で消える三次元ずつの部分の積は零で、補正後も指数成長する。
一つあればその二次数を$d,e$、対応射影直線上で$x_*,y_*$に対応する二つの異なる点を$P,Q$とする。
対応枝上の積は全$P_{d+e}$を張り、それと独立に水平枝のみ、垂直枝のみへの
$3-d,3-e$次元の像を持つ（対応枝への制限核と、相手での単元の積）。ここまでで七次元である。

さらに対応枝上の二切断を$f,g$と書くと、他の水平枝の一次係数の重み付き和から
$f'(P)g(P)$を、他の垂直枝から$f(Q)g'(Q)$を読み取れる。これは$x_*,y_*$での
一次の貼り合わせ式によるもので、自明化による零階項は対応枝の積の値である。
この二座標は上記の$3-d,3-e$次元の像では消える。対応枝の乗法核に制限すると、
二座標は対角の括弧で割った$(d-1,e-1)$形式の$(P,P),(Q,Q)$での評価に比例する。
非零の定数倍は各点の座標の選択による。従って$d+e\ge3$なら二条件は独立で$r\ge9$となり矛盾する。
$d=e=1$なら$r=8$であるが、対応枝への二つの制限核は各二次元で相互の積が零となる。
これも標準ASを排除する。

\paragraph{二つの特異点の対がグラフ上にある場合。}
各対応枝は根（それぞれの特異点の原像）を対応させる。
実際、$x_*$上に対応枝の根以外の$Y_j$の点があれば、その枝には別の垂直成分がないので、
この点は$y_*$と同じファイバー内で孤立してしまう。
また$y_i\ne y_*$の水平枝は、$x_*$のファイバーの連結性により
$x_*$上の垂直枝の根以外の点に付く。逆も同じである。
従って対応枝以外は中心$(x_*,y_*)$を通る水平または垂直な枝と、
その根以外に付く反対向きの末端枝だけである。付属枝を持たない中心枝自身も末端に数える。

必要な線型包の計算を与える。$e,f$を中心での両評価ベクトルとし、
各々一般の第一・第二切断で値が1となるように点評価を規格化する。
その切断の核を$U,V$とし、各枝で中心からの差を張る部分空間を$U_i,V_j$とする。
その次元は$d_i,e_j$である。$U_i$の直和から$U$への写像の唯一の核を$(h_i)_i$とし、
$V_j$では$(k_j)_j$とする。$h_i,k_j$は各々零でない根の接方向である。
これは各枝の値を一致させた上で一次係数に一式だけを課す表示による。
正規化に用いる切断は全ての付け根で非零、また各対応枝上で両側の零因子が交わらないように選べる。

点評価の包から$U\otimes V$の成分を捨て、中心の直線も法として$U\oplus V$に射影する。
各対応枝の中心を法とした$d_i+e_j$次元の包は$U_i\oplus V_j$へ同型になる。
実際、二つの切断を規格化切断で割った有理関数の空間は、極因子が互いに素なので定数だけを共有する。
従って$U_i\oplus0$のベクトル$h_i$をこの包に持ち上げたときの
$U_i\otimes V_j$成分$D_{ij}$が一意に決まる。対応枝の根での一次微分は混合成分が零で、
投影が$(h_i,c_{ij}k_j)$となる（$c_{ij}\ne0$）。従って$0\oplus k_j$の持ち上げの混合成分は
$-c_{ij}^{-1}D_{ij}$である。
$D_{ij}$は非零である。さらに$d_i,e_j$の少なくとも一つが2以上なら、
$D_{ij}$は$h_i\otimes k_j$の倍数ではない。
もしその倍数なら$h_i\otimes(f+z k_j)$という純テンソルが対応枝の線型包に入る。
積$P_d\otimes P_e\to P_{d+e}$の双対像の純テンソルは一つの点での二つの評価の積だけである。
実際、単項式座標の行列は$(a_{i+j})$で、階数1なら相隣る二行二列の行列式と全二次小行列式が消える。
非零な分解$u_i v_j$に対し$u_{i+1}v_j=u_i v_{j+1}$なので、
二列のシフトが非零なら両列は同じ公比の等比数列、端のシフトが零なら両方が同じ端の評価である。
これで純テンソルの主張が従う。次数2以上の有理正規曲線は根の接線と根以外で交わらない。
$h_i$は根の評価とは異なり、$f$だけを取ることも対応枝の同一パラメータに反する。
$d_i=e_j=1$なら$D_{ij}$は非零の$h_i\otimes k_j$倍で、対応する$f+z k_j$は
その直線上の根以外の点を表す。

水平枝の$U_i$方向の増分は$u\otimes(f+v_i)$であり、$f+v_i$はその枝の第二像の点である。
垂直枝の$V_j$方向は$(e+u_j)\otimes v$で同様である。従って全点評価の包は
中心と各枝のこれらの増分で張られる。付け根の評価が増分の基点として必要な場合も、
それは中心枝の線型包内にある。$U\oplus V$への投影は全射で、直和の添字空間での関係は
先の$(h_i)_i,(k_j)_j$の二つである。ゆえに$r=7+\dim\langle H,K\rangle$で、ここで
\[
 H=\sum_{(i,j)\ {\rm 対応}}D_{ij}+\sum_{i\ {\rm 水平}}h_i\otimes v_i,\qquad
 K=-\sum_{(i,j)\ {\rm 対応}}c_{ij}^{-1}D_{ij}+\sum_{j\ {\rm 垂直}}u_j\otimes k_j .
\]

対応枝がなければ$H=0$は全$v_i$が等しいことと同値である。
$h_i$の関係が全係数1の関係だけだからである。$K=0$も同様で、
$r=7$なら全水平枝が同じ点に、全垂直枝も同じ点に載る。両側でその点において消える
三次元ずつの切断の積が零となり、一次元の補正を許しても排除される。
対応枝があれば$r=8$となることも確認する。$d_i\ge2$または$e_j\ge2$の対応枝を除いた
連結な残りの包は、両射影の次元下界により少なくとも$9-d_i-e_j$次元である。
その枝だけで二重に消える側の切断（根で二重に消え、他枝上では零）と相手の全切断の積は、
残りで零となり、その枝で$P_{d_i+e_j-2}$の次元$d_i+e_j-1$を与える。
従って$r\ge8$である。全対応枝が二次数$(1,1)$なら、$H$は$\sum_i h_i\otimes v'_i$の形で、
対応枝の$f+v'_i$は対応する第二直線の根以外の点である。水平枝の第二像は根または
別の垂直枝上の点なのでこれと一致しない。他の対応枝の根以外の点も一致しない。
$p\ge2$より全$v'_i$は一致せず$H\ne0$である。この場合にも$r=8$である。

従って残りでは、実際の評価は$\Phi$を経由する。末端の水平枝の次数が2以上なら、
その枝だけに支えを持つ非零切断での評価は第二像の一点での値だけを通り、階数は1以下である。
また末端の水平枝が二つあれば、その二枝だけに支えを持つ非零切断がある
（根の値を零、一次係数の重み付き和を零とする）。その評価は高々二点の値を通る。
末端枝上の切断は、他の垂直枝が水平枝内の根以外に付いていないため、そのような垂直枝上に
追加の支持を持つことはない。従って末端水平枝は高々一本で次数1である。末端垂直枝も同様である。
付属枝を持つ中心の水平枝には高々一本の次数1の垂直枝が付く。この水平枝の次数が2以上なら、
その枝だけに支えを持つ第一切断を選べる。積の値は付属直線上の第二切断の全二座標で決まり、
水平枝上で使う定数もその直線の根での値である。評価が階数2以下となり矛盾する。
従って全ての非対応枝は次数1で、非対応の各群は一本の中心枝または一本ずつの水平・垂直枝の組である。

対応枝が一本もなければ各群が一つの末端を持つため高々二群、各群高々二枝である。
$p,q\ge2$なら$p=q=2$で各側二本の次数1の和が4に達せず矛盾する。
少なくとも一つ対応枝があるとする。$H,K$に残る非対応群の非零項は、根以外の付け根を持つ
水平・垂直枝の組についての非零の$h_i\otimes k_j$倍で、$H,K$の一方にだけ現れる。
これは両枝が直線のためである。これらの$(i,j)$と対応枝の$(i,j)$は添字を共有しない部分的な対応である。
このような対応の各項$D_{ij}$または$h_i\otimes k_j$は独立である。
実際、$D_{ij}$に$kh_i\otimes k_j$の外の成分があれば、左または右で
根の全接方向の包を法とした成分が非零になる。$U_i$と根の全接方向の包の交わりは$kh_i$であり、
それ以外の$U_i$の方向はその枝に独立だから、その係数は関係中で零である。右側も同様である。
残った$h_i\otimes k_j$では、左右どちらかが全添字の真部分ならその側の方向は独立である。
全単射かつ$p=q\ge3$なら最初の$p-1$方向を基底に取ると最後の積だけが非対角成分を持つ。
$p=q=2$で両項とも単純な接方向の積になるには双方の全枝が次数1でなければならず、
各側の総次数4に反する。独立性が示された。

$r\le8$のため$H,K$は従属しなければならない。対応項があるので両方非零であり、
片方にしか現れない付属枝の項は存在できない。従って非対応枝があれば全て中心を通る。
さらに対応枝全てで$c_{ij}$は同じスカラーである。
これにより、全枝上の関数の値を一致させ、一次係数の重み付き和を一式で課す
一つの楕円多重点曲線$Z$を構成できる。詳しくは各第一枝で$h_i$を、各第二枝で
$k_j$を定める接パラメータを使い、対応枝では同型の微分が共通の$c_{ij}$である。
第二側の全接パラメータを同じ定数で取り直せば、一方だけの枝も含めて同じ和零の条件に揃う。
$h_i,k_j$の規格化が構造層のこの条件を与えるのは、規格化切断で割った局所正則関数の
一次係数の唯一の関係を取ったためである。対応枝は同一とし、水平または垂直枝を一方で
縮める二射を得る。根での関数の条件が合うのでこれらはスキームの射である。

この$Z$上の埋込み束の引き戻しの切断は元の四座標全体で、
積束の切断は八次元である。後者は正規化での共通値と一次係数の一式から計算できる。
$\Phi$が階数8なのでこの積の切断を全て張り、実際の六式はその二次元鉛筆の表示となる。
基点は純テンソルを与えるから存在しない。命題\ref{prop:semipositive-pencil-balance}により
二つの多重次数は一致し、非対応枝は存在できない。従って$X,Y$は同じ$Z$であり、
多重点の全次数定理と同型復元が適用できる。
\end{proof}

\begin{corollary}[階数8の被約反対核に対する判定]\label{cor:all-reduced-apolar-recognition}
隣接階数6・反対階数8の両対称テンソルの二つの反対核が被約四次完全交叉を定めるとする。
標準ASとなる必要十分条件は、第一窓に純テンソル関係がなく、両曲線の同型のうち
二つの埋込み束の多重次数を一致させるもののグラフへの制限をテンソル汎関数が経由することである。
このとき同定と鉛筆が復元される。反対階数7の標準AS特殊化では、二つの極限鉛筆を
ともに被約四次完全交叉に保つことはできない。
\end{corollary}
\begin{proof}
全両対称AS窓の閉包定理と前定理で必要性を得る。十分性では同定した曲線上で
二つの束の積は全切断を張るので、階数6の隣接対合の核は二次元である。
この鉛筆の基点は純テンソル関係を与えるため存在しない。
各被約曲線型での全次数定理と復元定理を適用する。最後は前定理の階数の結論である。
\end{proof}
